%% =================================================================
%% Kingma, Ba. "Adam: A Method for Stochastic Optimization."
%% ICLR 2015. arXiv:1412.6980v9.
%%
%% Reconstruction of page 8 (printed p. 8) as study notes.
%% Generated by: reproduce-all-pages-basic-tex (batch 1-15)
%%
%% Structure: Figure 4 (bias-correction grid, 2x6 VAE loss plots)
%% + Sec 6.4 EXPERIMENT: BIAS-CORRECTION TERM + Sec 7 EXTENSIONS
%% + Sec 7.1 AdaMax with derivation equations (6), (7).
%% Equation counter: was 5, becomes 7.
%% =================================================================

\documentclass[11pt]{article}

\usepackage[margin=1in]{geometry}
\usepackage{amsmath, amssymb}
\usepackage{mathrsfs}
\usepackage{bm}
\usepackage{xcolor}
\usepackage{fancyhdr}

\pagestyle{fancy}
\fancyhf{}
\lhead{\small Published as a conference paper at ICLR 2015}
\rhead{}
\cfoot{\small 8 $\to$ \arabic{page}}
\renewcommand{\headrulewidth}{0.4pt}

\setcounter{page}{1}
\setcounter{equation}{5}

\begin{document}

% ---------- Figure 4 placeholder ----------
\noindent
\begin{center}
\fbox{\parbox{0.92\textwidth}{\centering\vspace{1em}
\textbf{[Figure 4 --- placeholder]}\\[0.8em]
\textbf{2$\times$6 grid} of VAE training-loss curves for a
parameter sweep over $\beta_1 \in \{0, 0.9\}$,
$\beta_2 \in \{0.99, 0.999, 0.9999\}$, and
$\log_{10}(\alpha) \in [-5, -1]$, comparing Adam \emph{with}
bias-correction terms (\textbf{red line}) vs.\ Adam
\emph{without} bias-correction terms (\textbf{green line}).
Y-axes plot loss ($\sim 100$ to $140$); x-axes plot
$\log_{10}(\alpha)$.\\[0.4em]
\textbf{(a) after 10 epochs (left three columns)}: at
$\beta_2 = 0.99$ the two curves are nearly coincident; at
$\beta_2 = 0.999$ and especially $\beta_2 = 0.9999$, the green
(no-correction) curve lies \emph{above} the red (with-correction)
curve --- bias correction reduces loss. The gap grows as
$\beta_2 \to 1$.\\[0.4em]
\textbf{(b) after 100 epochs (right three columns)}: same pattern
but the gap is larger and more systematic. Without bias
correction and with $\beta_2$ close to $1$, training is visibly
worse --- the green curve rises at small $\alpha$ values where
the red curve stays flat, indicating the no-bias-correction
version is unstable in the regime that's actually needed for
sparse gradients.\\[0.4em]
\textbf{Takeaway}: bias correction matters most exactly when
$\beta_2 \to 1$, which is the regime required for sparse
gradients. Removing the correction turns Adam into a version of
momentum-RMSProp that is numerically unstable in this useful
corner of parameter space.
\vspace{1em}}}
\end{center}
\vspace{0.3em}
\begin{center}
\textbf{Figure 4:} Effect of bias-correction terms (red line)
versus no bias correction terms (green line) after 10 epochs
(left) and 100 epochs (right) on the loss (y-axes) when learning
a Variational Auto-Encoder (VAE) (Kingma \& Welling, 2013), for
different settings of stepsize $\alpha$ (x-axes) and
hyper-parameters $\beta_1$ and $\beta_2$.
\end{center}

\vspace{1em}

% ---------- Section 6.4 ----------
\subsection*{6.4 \textsc{Experiment: Bias-correction term}}

We also empirically evaluate the effect of the bias correction
terms explained in sections~2 and~3. Discussed in section~5,
removal of the bias correction terms results in a version of
RMSProp (Tieleman \& Hinton, 2012) with momentum. We vary the
$\beta_1$ and $\beta_2$ when training a variational auto-encoder
(VAE) (Kingma \& Welling, 2013) with the same architecture as in
(Kingma \& Welling, 2013) with a single hidden layer with 500
hidden units with softplus nonlinearities and a 50-dimensional
spherical Gaussian latent variable. We iterated over a broad
range of hyper-parameter choices, i.e.\
$\beta_1 \in [0, 0.9]$ and
$\beta_2 \in [0.99, 0.999, 0.9999]$, and
$\log_{10}(\alpha) \in [-5, \ldots, -1]$. Values of $\beta_2$
close to $1$, required for robustness to sparse gradients,
results in larger initialization bias; therefore we expect the
bias correction term is important in such cases of slow decay,
preventing an adverse effect on optimization.

In Figure~4, values $\beta_2$ close to $1$ indeed lead to
instabilities in training when no bias correction term was
present, especially at first few epochs of the training. The
best results were achieved with small values of $(1 - \beta_2)$
and bias correction; this was more apparent towards the end of
optimization when gradients tends to become sparser as hidden
units specialize to specific patterns. In summary, Adam performed
equal or better than RMSProp, regardless of hyper-parameter
setting.

\vspace{0.5em}

\section*{7 \textsc{Extensions}}

\subsection*{7.1 \textsc{AdaMax}}

In Adam, the update rule for individual weights is to scale
their gradients inversely proportional to a (scaled) $L^2$ norm
of their individual current and past gradients. We can
generalize the $L^2$ norm based update rule to a $L^p$ norm
based update rule. Such variants become numerically unstable for
large $p$. However, in the special case where we let
$p \to \infty$, a surprisingly simple and stable algorithm
emerges; see algorithm~2. We'll now derive the algorithm. Let,
in case of the $L^p$ norm, the stepsize at time $t$ be inversely
proportional to $v_t^{1/p}$, where:
%
\begin{align}
v_t &= \beta_2^p \, v_{t-1} + (1 - \beta_2^p) |g_t|^p \\
    &= (1 - \beta_2^p) \sum_{i=1}^{t} \beta_2^{p(t-i)} \cdot |g_i|^p
\end{align}

% [page ends here — continues on page 9]

\end{document}
